\scanpage{45}
On aura dans $V$
\[
(4)\qquad\sum_{j\in J}R_j=0
\]
ce qui donne un isomorphisme
\[
(5)\qquad V\simeq V_J(\mathbb R)=\operatorname{Ker}(\mathbb R^J\xrightarrow[\text{somme}]{}\mathbb R)
\]
et on a, par la $3^{\unsure{\text{ième}|\text{ème}}}$ formule (3)
\[
(6)\qquad\Delta=V^\perp\simeq\mathbb R(\varepsilon)
\]
\margin{``axe polair''}
\margin{orth. de $V$ dans $E$}
\margin{(isom compatible avec la structure quadratique)}
donc
\[
(7)\qquad E\simeq V_J(\mathbb R)\oplus\mathbb R(\varepsilon)
\]
\margin{isom d'espaces euclidiens}
où dans $V_J(\mathbb R)$ on prend la seule forme quadratique \inserted{$q$} pour laquelle les $R_j$ \deleted{sont} \inserted{satisfont} \deleted{de carré \unsure{un}}
\[
q(R_j)=1
\]
\margin{Il est clair, par cette caractérisation de $q$, que l'opération naturelle de $\unsure{\mathcal G_J|G_J}$ sur $V(J)$ conserve $q$.}
\unsure{(prenant|(prendre} une base formée de $R_i,R_j$, $R_{\unsure{n|h}}$ de coordonnées $(-1,-1)$, la forme quadratique en question est $x^2-xy+y^2$)

Dans le plan euclidien $V$, les $R_i$ \figure{\begin{tikzpicture}[baseline=0pt]\draw[->] (.3,.55)--(0,.35)--(0,0);\end{tikzpicture}} forment les sommets d'un triangle équilatéral inscrit dans le cercle unité. L'opération de $\unsure{\mathcal G_J|G_J}$ sur $E$ se fait par son opération \deleted{sur} tautologique sur $V(J)$ et l'opération par $\operatorname{sg}(g)$ sur $\Delta$:
\margin{44}
